\documentclass[10pt,a4paper]{amsart}
\usepackage[margin=19mm]{geometry}
\usepackage{amsmath,amssymb,mathtools}
\usepackage[scr=rsfs,cal=euler]{mathalfa}
\usepackage{xfrac}
\usepackage[unicode,hidelinks,bookmarks=false]{hyperref}
\newcommand{\ie}{i.e.\ }
\newcommand{\cf}{cf.\ }
\newcommand{\resp}{respectively\ }
\newcommand{\Subsection}{\subsection}
\providecommand{\nobreakdash}{\nobreak\hbox{-}}
\setlength{\emergencystretch}{2em}
\multlinegap0pt
\usepackage{array,amsxtra,hhline,longtable,enumitem}
\newcolumntype{P}[1]{>{\raggedright\let\newline\\\arraybackslash\hspace{0pt}}m{#1}}
\providecommand{\C}{\mathcal C}
\providecommand{\U}{\mathcal U}
\newtheorem{theorem}{Theorem}[section]
\newtheorem{thm}[theorem]{Theorem}
\newtheorem{lemma}[theorem]{Lemma}
\newtheorem{prop}[theorem]{Proposition}
\newtheorem{coro}[theorem]{Corollary}
\theoremstyle{definition}
\newtheorem{defi}[theorem]{Definition}
\newtheorem{rema}[theorem]{Remark}
\DeclareMathOperator{\Ram}{Ram}
\DeclareMathOperator{\uf}{uf}
\DeclareMathOperator{\fr}{fr}
\DeclareMathOperator{\cov}{cov}
\DeclareMathOperator{\Cone}{Cone}
\DeclareMathOperator{\Star}{Star}
\DeclareMathOperator{\Lin}{Lin}
\DeclareMathOperator{\Span}{Span}
\DeclareMathOperator{\re}{re}
\DeclareMathOperator{\im}{im}
\DeclareMathOperator{\Hom}{Hom}
\DeclareMathOperator{\Supp}{Supp}
\DeclareMathOperator{\Scat}{Scat}
\DeclareMathOperator{\Fan}{Fan}
\DeclareMathOperator{\ChamberFan}{ChamberFan}
\DeclareMathOperator{\Clear}{Clear}
\DeclareMathOperator{\Nar}{Nar}
\DeclareMathOperator{\ScatFan}{ScatFan}
\DeclareMathOperator{\CambScat}{CambScat}
\newcommand{\gFan}{\g\!\Fan}

\newcommand{\newword}[1]{\emph{#1}}


\newcommand\llbracket{[\![}
\newcommand\rrbracket{]\!]}

\newcommand{\integers}{\mathbb Z}
\newcommand{\rationals}{\mathbb Q}
\newcommand{\reals}{\mathbb R}
\newcommand{\ep}{\varepsilon}
\newcommand{\thet}{\vartheta}
\newcommand{\set}[1]{{\left\lbrace #1 \right\rbrace}}
\newcommand{\br}[1]{{\langle #1 \rangle}}
\newcommand{\A}{{\mathcal A}}
\newcommand{\F}{{\mathcal F}}
\newcommand{\D}{{\mathfrak D}}
\newcommand{\N}{{\mathcal N}}
\newcommand{\p}{{\mathfrak p}}
\newcommand{\ck}{\spcheck}

\renewcommand{\th}{^\text{th}}
\newcommand{\g}{\mathbf{g}}
\newcommand{\s}{\mathbf{s}}
\newcommand{\m}{\mathbf{m}}
\newcommand{\cc}{\mathbf{c}}
\renewcommand{\b}{\mathbf{b}}
\renewcommand{\k}{\Bbbk}
\renewcommand{\a}{\mathbf{a}}
\renewcommand{\t}{\mathbf{t}}
\renewcommand{\v}{\mathbf{v}}
\renewcommand{\u}{\mathbf{u}}
\newcommand{\tB}{\tilde{B}}
\renewcommand{\C}{\mathcal{C}}
\newcommand{\M}{\mathcal{M}}
\renewcommand{\U}{\mathcal{U}}
\renewcommand{\O}{\mathcal{O}}
\renewcommand{\H}{\mathcal{H}}
\renewcommand{\P}{\mathbb{P}}
\renewcommand{\S}{\mathbf{S}}
\renewcommand{\d}{{\mathfrak d}}
\newcommand{\seg}[1]{\overline{#1}}
\newcommand{\hy}{\hat{y}}


%%%%%%%%% a placer avant \begin{document}


%%%%%%%%%%%%%%%%%%%%%%%%%%%%%%%%%%%%%
\newcommand*{\mk}{\mkern -1mu}
\newcommand*{\Mk}{\mkern -2mu}
\newcommand*{\mK}{\mkern 1mu}
\newcommand*{\MK}{\mkern 2mu}

\hypersetup{urlcolor=purple, linkcolor=blue, citecolor=red}


\newcommand*{\romanenumi}{\renewcommand*{\theenumi}{\roman{enumi}}}
\newcommand*{\Romanenumi}{\renewcommand*{\theenumi}{\Roman{enumi}}}
\newcommand*{\alphenumi}{\renewcommand*{\theenumi}{\alph{enumi}}}
\newcommand*{\Alphenumi}{\renewcommand*{\theenumi}{\Alph{enumi}}}
%%%%%%%%%%%%%%%%%%%%%%%%%%%%%%%%%%


\title[Affine limiting scattering wall]{A combinatorial approach to scattering diagrams: original Theorem3.4}
\author{Nathan Reading}
\date{Source: Algebraic Combinatorics3(3)(2020),603--636; DOI10.5802/alco.107}
\begin{document}\maketitle
\noindent\textit{Editorial scope.} The counted unit is original Theorem3.4. The complete root, coefficient, scattering and path-limit definitions and both original coefficient-division proofs are retained as uncounted core. Finite root illustrations and rank-two example pictures are omitted; all their data used in the argument are explicitly defined in the retained text. Original equation(21) has a square missing from the term following $-16$, whereas the preceding original equation(20) explicitly contains the product of the two corresponding sums; the printed formula and complete author argument remain verbatim.
\setcounter{section}{1}
\section{Transposed cluster scattering diagrams with principal coefficients}\label{transp sec}
In this section, we begin with an exchange matrix and construct a cluster scattering diagram with principal coefficients.
We introduce a global transpose in order to make a cluster monomial $\thet_{m_0}$ have $\g$-vector $m_0$, in the tall-extended-exchange-matrix sense of~\cite{ca4} (as discussed in the Introduction).
A useful side effect of the transpose is that in acyclic finite type the scattering fan for an exchange matrix coincides with the Cambrian fan.
The connection to the Cambrian fan will be made in Section~\href{https://alco.centre-mersenne.org/articles/10.5802/alco.107/}{4}.
We also use root and weight lattices as part of our initial data.
%SinceArXiv: Changed this to cite Bridgeland and GHKK more equally
Some motivation for this choice is provided by a theorem of~\cite{Bridgeland,GHKK}, quoted below as Theorem~\ref{root thm}, and by the Cambrian fan construction in acyclic finite type.
Except for the transpose and the language of root and weight lattices, our implementation of principal coefficients follows~\cite[Appendix~B]{GHKK}.

\subsection{Exchange matrices, root systems and Coxeter groups}\label{back sec}
We start with an \newword{exchange matrix} $B=[b_{ij}]$, a square integer matrix indexed by $\set{1,\ldots,n}$ that is \newword{skew-symmetrizable}, meaning that there exist real numbers $\delta_i$ such that $\delta_ib_{ij}=-\delta_jb_{ji}$ for all $i,j\in\set{1,\ldots,n}$.
We choose the $\delta_i$ so that $\delta_i^{-1}$ is an integer for each $i$ and $\gcd(\delta_i^{-1}:i\in\set{1,\ldots,n})=1$.
We can do this because $B$ is an integer matrix.
We follow the usual convention and call $n$ the \newword{rank} of $B$.
This conflicts with the usual definition of rank, but not badly:
The constructions considered here require that $B$ be extended, by adjoining rows or columns as discussed above, to obtain a matrix of rank $n$ in the usual sense.

Let $A$ be the \newword{Cartan matrix} associated to $B$.
That is, $A=[a_{ij}]$ where $a_{ii}=2$ for $i=1,\ldots,n$ and $a_{ij}=-|b_{ij}|$ for all distinct $i,j\in\set{1,\ldots,n}$.
In particular, $A$ is \newword{symmetrizable} because $\delta_ia_{ij}=\delta_ja_{ji}$ for all $i,j\in\set{1,\ldots,n}$.

Choose a real vector space $V$ with a distinguished basis $\alpha_1,\ldots,\alpha_n$ called the \newword{simple roots}.
The lattice $Q=\Span_\integers(\alpha_1,\ldots,\alpha_n)$ is called the \newword{root lattice}.
Define the \newword{simple co-roots} to be $\alpha_i\ck=\delta_i^{-1}\alpha_i$, so that $\alpha_1\ck,\ldots,\alpha_n\ck$ is another basis for $V$.
The lattice $Q\ck=\Span_\integers(\alpha_1\ck,\ldots,\alpha_n\ck)$ is the \newword{co-root lattice}.
Since the $\delta_i$ were chosen so that each $\delta_i^{-1}$ is an integer, $Q\ck$ is a sublattice of $Q$ of finite index.

Given a primitive vector $\beta$ in $Q$ (an element $\beta\in Q$ not equal to $k\beta'$ for $k>1$ and $\beta'\in Q$), write $\beta\ck$ for the primitive vector in $Q\ck$ that is a positive scaling of $\beta$.
Given primitive $\beta\ck\in Q\ck$, write $\beta$ for the corresponding primitive vector in $Q$.

Let $K:V\times V\to\reals$ be the bilinear form defined, in the basis of simple roots on the right and simple co-roots on the left, by $K(\alpha_i\ck,\alpha_j)=a_{ij}$.
This restricts to an integer-valued form $K:Q\ck\times Q\to\integers$.
The form is symmetric because 
\[K(\alpha_i,\alpha_j)=\delta_iK(\alpha_i\ck,\alpha_j)=\delta_ia_{ij}=\delta_ja_{ji}=\delta_jK(\alpha_j\ck,\alpha_i)=K(\alpha_j,\alpha_i).\]
Let $\omega:V\times V\to\reals$ be the bilinear form defined by $\omega(\alpha_i\ck,\alpha_j)=b_{ij}$.
This takes integer values on $Q\ck\times Q$ and is skew-symmetric by a similar calculation.

Let $V^*$ be the dual vector space to $V$ and let $\br{\,\cdot\,,\,\cdot\,}:V^*\times V\to\reals$ be the usual pairing.
Define the \newword{fundamental weights} to be the basis $\rho_1,\ldots,\rho_n$ for $V^*$ that is dual to $\alpha_1\ck,\ldots,\alpha_n\ck$, in the sense that $[\br{\rho_i,\alpha\ck_j}]$ is the identity matrix.
(The fundamental weights are dual to the simple \emph{co}-roots, not the simple roots.)
The lattice $P=\Span_\integers(\rho_1,\ldots,\rho_n)$ is called the \newword{weight lattice}.
Define the \newword{fundamental co-weights} to be the basis $\rho_1\ck,\ldots,\rho_n\ck$ for $V^*$ that is dual to $\alpha_1,\ldots,\alpha_n$.
We have $\rho_i=\delta_i\rho_i\ck$ for all $i$.
%(Since $\alpha\ck_i=\delta_i^{-1}\alpha_i$ for each $i\in\set{1,\ldots,n}$, the matrix $[\br{\rho_i,\delta_j^{-1}\alpha_j}]$ is the identity, so $[\br{\delta_j^{-1}\rho_i,\alpha_j}]$ is the identity,
%Thus $\rho_i\ck=\delta_i^{-1}\rho_i$.)
The lattice $P\ck=\Span_\integers(\rho_1\ck,\ldots,\rho_n\ck)$ is the \newword{co-weight lattice}.
Since each $\delta_i^{-1}$ is an integer, $P$ is a superlattice of $P\ck$ of finite index.

The \newword{dominant chamber} in $V^*$ is the full-dimensional simplicial cone
\begin{equation}\label{D def}
D=\bigcap_{i=1}^n \set{p\in V^*: \br{p,\alpha_i}\ge 0}.
\end{equation}
Equivalently, $D$ is the nonnegative real span of the fundamental weights or of the fundamental co-weights.




For each $i=1,\ldots,n$, define $s_i$ to be the reflection on $V$ given by $s_i(\alpha_j)=\alpha_j-a_{ij}\alpha_i$.
The set $S$ of \newword{simple reflections} $s_i$ generates a \newword{Coxeter group} $W$.
(More precisely, $(W,S)$ is a Coxeter system.)
%This means that $W$ has a presentation with generators $S=\set{s_1,\ldots,s_n}$ and a certain form of relations.)
The action of $s_i$ on $V$ defines an action on $V^*$ in the usual way.
Namely, $s_i$ sends $\lambda\in V^*$ to the unique vector $s_i\lambda\in V^*$ such that $\br{s_i\lambda,s_i\beta}=\br{\lambda,\beta}$ for all $\beta\in V$.


The \newword{real roots} $\Phi^{\re}$ associated to $A$ are the vectors of the form $w\alpha_i$ for $w\in W$ and $i=1,\ldots,n$.
When $\Phi^{\re}$ is infinite, there are also \newword{imaginary} roots $\Phi^{\im}$ associated to $A$, which we need not define here.
The \newword{root system} associated to $A$ is the disjoint union $\Phi=\Phi^{\re}\cup\Phi^{\im}$.
(Root systems as we have defined them are sometimes called \newword{crystallographic}; there is a more general notion that does not concern us here.)
The root system $\Phi$ is also a disjoint union $\Phi=\Phi_+\cup\Phi_-$ such that each $\beta\in\Phi_+$ is a nonnegative linear combination of simple roots and each $\beta\in\Phi_-$ is a nonpositive linear combination of simple roots.
The roots in $\Phi_+$ are called the \newword{positive roots}.
We write $Q^+$ for the subset of $Q$ consisting of nonzero vectors obtained as nonnegative integer combinations of simple roots. 
%We similarly define the \newword{dual root system} $\Phi\ck=\Phi\ck_+\cup\Phi\ck_-$ and write $Q\ck_+$ for the subset of $Q\ck$ consisting of nonnegative integer combinations of simple co-roots.


The \newword{real co-roots} %$\Phi\ck$
are the vectors of the form $w\alpha\ck_i$ for $w\in W$ and $i=1,\ldots,n$.
The vector $\beta\ck$ is a real co-root if and only if its scaling $\beta$ is a real root.

\setcounter{theorem}{1}
\subsection{Principal-coefficients transposed scattering diagrams in root notation}\label{prin trans sec}
Table~\ref{scat root data} describes the initial data for a transposed scattering diagram with principal coefficients.
The only input is an exchange matrix $B=[b_{ij}]$, from which we extract a Cartan matrix and make the definitions of Section~\ref{back sec}.
%\begin{table}[ht]
%\begin{tabular}{P{70pt}|P{260pt}}
%Notation&Description/requirements\\\hline\\[-6pt]
%%$B=[b_{ij}]_{1\le i,j\le n}$&skew-symmetrizable exchange matrix\\\hline
%$Q^+$&$\set{\sum_{i=1,\ldots,n}a_i\alpha_i:a_i\in\integers,\,a_i\ge0,\,\sum_{i\in I_{\uf}}a_i>0}$\newline positive part of root lattice \\\hline
%%$p^*:Q\to P\oplus Q$&$p^*(\alpha_i)=(\sum_{j=1}^nb_{ji}\rho_i,\,\alpha_i)$\\\hline
%$x_1,\ldots,x_n$\newline $y_1,\ldots,y_n$&indeterminates\\\hline
%$x^\lambda y^\beta$&$x_1^{a_1}\cdots x_n^{a_n}y_1^{c_1}\cdots y_n^{c_n}$ for $(\lambda,\beta)=(\sum_{i=1}^na_i\rho_i,\sum_{i=1}^nc_i\alpha_i)\in P\oplus Q$\\\hline
%$\hy_1,\ldots,\hy_n$&$\hy_i=y_ix_1^{b_{1i}}\cdots x_n^{b_{ni}}$\\\hline
%$x^\lambda \hy^\beta$&
%$x_1^{a_1}\cdots x_n^{a_n}\hy_1^{c_1}\cdots\hy_n^{c_n}$ for $(\lambda,\beta)=(\sum_{i=1}^na_i\rho_i,\sum_{i=1}^nc_i\alpha_i)\in P\oplus Q$\\\hline
%$\k$&a field of characteristic zero\\\hline
%$\k\llbracket \hy\rrbracket $&$\k\llbracket \hy_1,\ldots,\hy_n\rrbracket $\\\hline
%$\k\llbracket x,\hy\rrbracket $&$\k\llbracket x_1,\ldots,x_n,\hy_1,\ldots,\hy_n\rrbracket $\\\hline
%$\m$&ideal in $\k\llbracket x,\hy\rrbracket $ consisting of series with constant term zero\\\hline
%\end{tabular}
%\caption{Initial data for a transposed scattering diagram with principal coefficients}\label{scat root data}
%\end{table}

{\renewcommand{\arraystretch}{1.2}
\begin{longtable}{P{70pt}|P{260pt}}
\caption{Initial data for a transposed scattering diagram with principal coefficients}\label{scat root data}\\
Notation&Description/requirements\\\hline\hline
\endfirsthead
\caption{(continued)}\\
Notation&Description/requirements\\\hline\hline
\endhead
%$B=[b_{ij}]_{1\le i,j\le n}$&skew-symmetrizable exchange matrix\\\hline
$Q^+$&$\set{\sum_{i=1,\ldots,n}a_i\alpha_i:a_i\in\integers,\,a_i\ge0,\,\sum_{i=1,\ldots,n}a_i>0}$\vrule height14pt depth0pt width0pt\newline positive part of root lattice \\\hline
%$p^*:Q\to P\oplus Q$&$p^*(\alpha_i)=(\sum_{j=1}^nb_{ji}\rho_i,\,\alpha_i)$\\\hline
$x_1,\ldots,x_n$\newline $y_1,\ldots,y_n$&indeterminates\\\hline
$x^\lambda y^\beta$&$x_1^{a_1}\cdots x_n^{a_n}y_1^{c_1}\cdots y_n^{c_n}$ for $(\lambda,\beta)=(\sum_{i=1}^na_i\rho_i,\sum_{i=1}^nc_i\alpha_i)\in P\oplus Q$\\\hline
$\hy_1,\ldots,\hy_n$&$\hy_i=y_ix_1^{b_{1i}}\cdots x_n^{b_{ni}}$\\\hline
$x^\lambda \hy^\beta$&
$x_1^{a_1}\cdots x_n^{a_n}\hy_1^{c_1}\cdots\hy_n^{c_n}$ for $(\lambda,\beta)=(\sum_{i=1}^na_i\rho_i,\sum_{i=1}^nc_i\alpha_i)\in P\oplus Q$\\\hline
$\k$&a field of characteristic zero\\\hline
$\k\llbracket \hy\rrbracket $&$\k\llbracket \hy_1,\ldots,\hy_n\rrbracket $\\\hline
$\k\llbracket x,\hy\rrbracket $&$\k\llbracket x_1,\ldots,x_n,\hy_1,\ldots,\hy_n\rrbracket $\\\hline
$\m$&ideal in $\k\llbracket x,\hy\rrbracket $ consisting of series with constant term zero\\\hline
\end{longtable}}%

As discussed in the Introduction, we work with a global transpose, relative to~\cite{GHKK}.
To avoid confusion, we will be explicit about this transpose in terminology and notation.
We also follow~\cite{scatfan} in working in a lower-dimensional space than~\cite{GHKK}.
For details, including an explanation of why the additional dimensions are unnecessary, see~\cite[Remarks~2.1, 2.12,~2.13]{scatfan}.

For the purpose of comparison, Table~\ref{scat root dict} gives this initial data in the context of the more general setup of~\cite{GHKK}.
Table~\ref{scat root dict} is designed for easy comparison with~\cite[Table~1]{scatfan}. 
The general setup leaves several choices, and we have made these choices in ways that are natural to the root-system context. 
We see the explicit transpose in the line of Table~\ref{scat root dict} describing $\epsilon_{ij}$.
%\renewcommand*{\arraystretch}{1.5}
%\begin{table}
%\begin{tabular}{P{88pt}|P{242pt}}
%General setup&Transposed, principal coefficients\\\hline
%$N$&$Q\oplus P$\quad (root lattice and weight lattice)\\\hline
%$M=\Hom(N,\integers)$&$P\ck\oplus Q\ck$\quad (co-weight lattice and co-root lattice)\\\hline
%$N_{\uf}\subseteq N$&Root lattice $Q$, identified with $Q\oplus\set{0}\subset Q\oplus P$\\\hline
%$I$&Two copies of $\set{1,\ldots,n}$\\\hline
%$I_{\uf}\subseteq I$&$\set{1,\ldots,n}$\\\hline
%$\set{\,\cdot\,,\,\cdot\,}:N\times N\to\rationals$&$\set{(\beta_1,\lambda_1),(\beta_2,\lambda_2)}=-\omega(\beta_1,\beta_2)+\br{\lambda_2,\beta_1}-\br{\lambda_1,\beta_2}$\\\hline
%$N^\circ\subseteq N$&$Q\ck\oplus P\ck$\quad (co-root lattice and co-weight lattice)\\\hline
%$M^\circ=\Hom(N^\circ,\integers)$&$P\oplus Q$\quad (weight lattice and root lattice)\\\hline
%$d_i$&$d_i=\delta_i^{-1}=\frac{K(\alpha_i,\alpha_i)}2=\frac2{K(\alpha_i\ck,\alpha_i\ck)}$ \quad so $\alpha_i\ck=d_i\alpha_i$\newline
%$d_i$ same on both copies of $\set{1,\ldots,n}$, so $\rho_i\ck=d_i\rho_i$\\\hline
%$\s=(e_i:i\in I)$&$\s=(\alpha_1,\ldots,\alpha_n,\rho_1,\ldots,\rho_n)$\quad basis for $Q\oplus P$\\\hline
%$N^+=N^+_\s$&$Q^+=\set{\sum_{i=1,\ldots,n}a_i\alpha_i:a_i\in\integers,\,a_i\ge0,\,\sum_{i\in I_{\uf}}a_i>0}$\newline positive part of root lattice \\\hline
%$[\,\cdot\,,\,\cdot\,]_\s:N\times N\to\rationals$&
%$[\alpha_i,\alpha_j]_\s=\set{\alpha_i,\alpha_j\ck}=\omega(\alpha_j\ck,\alpha_i)=b_{ji}$\newline
%$[\alpha_i,\rho_j]_\s=\br{\rho_j\ck,\alpha_i}=\delta_{ij}$ (Kronecker delta)\newline
%$[\rho_i,\alpha_j]_\s=\br{\rho_i,\alpha_j\ck}=-\delta_{ij}$ \newline
%$[\rho_i,\rho_j]_\s=0$
%\\\hline
%$\epsilon_{ij}=[e_i,e_j]_\s$&entries in matrix \raisebox{0pt}[14pt][9pt]{$\begin{bsmallmatrix}\,B^T&\,I\\-I&0\end{bsmallmatrix}$} ($n\times n$ blocks)\\\hline
%$(e^*_i:i\in I)$&$(\rho_i\ck,\ldots,\rho_n\ck,\alpha_1\ck,\ldots,\alpha_n\ck)$\quad basis for $P\ck\oplus Q\ck$\\\hline
%$(f_i:i\in I)$&$(\rho_i,\ldots,\rho_n,\alpha_1,\ldots,\alpha_n)$\quad basis for $P\oplus Q$\\\hline
%$p^*:N_{\uf}\to M^\circ$&$p^*(\alpha_i)=(\sum_{j=1}^nb_{ji}\rho_j,\,\alpha_i)$ %,\quad $p^*(\rho_i)=-\rho_i$
%\newline
%$p^*(\beta)=(\phi\ck\ni Q\ck\mapsto\omega(\phi\ck,\beta),\beta)$ for $\beta\in Q$
%\\\hline
%$(v_i\in M^\circ:i\in I)$&$(p^*(\alpha_1),\ldots,p^*(\alpha_n),\rho_1,\ldots,\rho_n)$\\\hline
%$(z_i:i\in I)$&$(x_1,\ldots,x_n,y_1,\ldots,y_n)$\quad indeterminates\\\hline
%$z^{(\lambda,\beta)}$&$x^\lambda y^\beta=x_1^{a_1}\cdots x_n^{a_n}y_1^{c_1}\cdots y_n^{c_n}$ for $(\lambda,\beta)=(\sum_{i=1}^na_i\rho_i,\sum_{i=1}^nc_i\alpha_i)\in P\oplus Q$\\\hline
%$(\zeta_i=z^{v_i}:i\in I)$&$(\hy_1,\ldots,\hy_n,x_1,\ldots,x_n)$ for $\hy_i=y_ix_1^{b_{1i}}\cdots x_n^{b_{ni}}$\\\hline
%$\zeta^{(\beta,\lambda)}$&
%$x^\lambda \hy^\beta=x_1^{a_1}\cdots x_n^{a_n}\hy_1^{c_1}\cdots\hy_n^{c_n}$ for $(\beta,\lambda)=(\sum_{i=1}^nc_i\alpha_i,\sum_{i=1}^na_i\rho_i)\in Q\oplus P$\\\hline
%$\k$&a field of characteristic zero\\\hline
%$\k\llbracket \zeta\rrbracket $&$\k\llbracket x,\hy\rrbracket =\k\llbracket x_1,\ldots,x_n,\hy_1,\ldots,\hy_n\rrbracket $\\\hline
%$\m$&ideal in $\k\llbracket x,\hy\rrbracket $ consisting of series with constant term zero\\\hline
%\end{tabular} 
%\caption{Initial data for a transposed scattering diagram with principal coefficients, in the context of the general setup}\label{scat root dict}
%\end{table} 

{\renewcommand{\arraystretch}{1.2}
\begin{longtable}{P{88pt}|P{242pt}}
\caption{Initial data for a transposed scattering diagram with principal coefficients, in the context of the general setup}\label{scat root dict}\\
General setup&Transposed, principal coefficients\\\hline\hline
\endfirsthead
\caption{(continued)}\\
General setup&Transposed, principal coefficients\\\hline\hline
\endhead
$N$&$Q\oplus P$\quad (root lattice and weight lattice)\\\hline
$M=\Hom(N,\integers)$&$P\ck\oplus Q\ck$\quad (co-weight lattice and co-root lattice)\\\hline
$N_{\uf}\subseteq N$&Root lattice $Q$, identified with $Q\oplus\set{0}\subset Q\oplus P$\\\hline
$I$&Two copies of $\set{1,\ldots,n}$, one indexing the simple roots in $N$ and one indexing the fundamental weights in $N$\\\hline
$I_{\uf}\subseteq I$&$\set{1,\ldots,n}$, indexing the simple roots in $N$\\\hline
$\set{\,\cdot\,,\,\cdot\,}:N\times N\to\rationals$&$\set{(\beta_1,\lambda_1),(\beta_2,\lambda_2)}=-\omega(\beta_1,\beta_2)+\br{\lambda_2,\beta_1}-\br{\lambda_1,\beta_2}$\\\hline
$N^\circ\subseteq N$&$Q\ck\oplus P\ck$\quad (co-root lattice and co-weight lattice)\\\hline
$M^\circ=\Hom(N^\circ,\integers)$&$P\oplus Q$\quad (weight lattice and root lattice)\\\hline
$d_i$&$d_i=\delta_i^{-1}=\frac{K(\alpha_i,\alpha_i)}2=\frac2{K(\alpha_i\ck,\alpha_i\ck)}$ \quad so $\alpha_i\ck=d_i\alpha_i$\newline
$d_i$ same on both copies of $\set{1,\ldots,n}$, so $\rho_i\ck=d_i\rho_i$\\\hline
$\s=(e_i:i\in I)$&$\s=(\alpha_1,\ldots,\alpha_n,\rho_1,\ldots,\rho_n)$\quad basis for $Q\oplus P$\\\hline
$N^+=N^+_\s$&$Q^+=\set{\sum_{i=1}^na_i\alpha_i:a_i\in\integers,\,a_i\ge0,\,\sum_{i=1}^na_i>0}$\newline positive part of root lattice \\\hline
$[\,\cdot\,,\,\cdot\,]_\s:N\times N\to\rationals$&
$[\alpha_i,\alpha_j]_\s=\set{\alpha_i,\alpha_j\ck}=\omega(\alpha_j\ck,\alpha_i)=b_{ji}$\newline
$[\alpha_i,\rho_j]_\s=\br{\rho_j\ck,\alpha_i}=\delta_{ij}$ (Kronecker delta)\newline
%SinceALCO: Fixed following (added ``-'' in front of \br{\rho_i,\alpha_j\ck})
$[\rho_i,\alpha_j]_\s=-\br{\rho_i,\alpha_j\ck}=-\delta_{ij}$ \newline
$[\rho_i,\rho_j]_\s=0$
\\\hline
$\epsilon_{ij}=[e_i,e_j]_\s$&entries in matrix \raisebox{0pt}[14pt][9pt]{$\begin{bsmallmatrix}\,B^T&\,I\\-I&0\end{bsmallmatrix}$} ($n\times n$ blocks)\\\hline
$(e^*_i:i\in I)$&$(\rho_i\ck,\ldots,\rho_n\ck,\alpha_1\ck,\ldots,\alpha_n\ck)$\quad basis for $P\ck\oplus Q\ck$\\\hline
$(f_i:i\in I)$&$(\rho_i,\ldots,\rho_n,\alpha_1,\ldots,\alpha_n)$\quad basis for $P\oplus Q$\\\hline
$p^*:N_{\uf}\to M^\circ$&$p^*(\alpha_i)=(\sum_{j=1}^nb_{ji}\rho_j,\,\alpha_i)$ %,\quad $p^*(\rho_i)=-\rho_i$
\newline
$p^*(\beta)=(\phi\ck\ni Q\ck\mapsto\omega(\phi\ck,\beta),\beta)$ for $\beta\in Q$
\\\hline
$(v_i\in M^\circ:i\in I)$&$(p^*(\alpha_1),\ldots,p^*(\alpha_n),\rho_1,\ldots,\rho_n)$\\\hline
$(z_i:i\in I)$&$(x_1,\ldots,x_n,y_1,\ldots,y_n)$\quad indeterminates\\\hline
$z^{(\lambda,\beta)}$&$x^\lambda y^\beta=x_1^{a_1}\cdots x_n^{a_n}y_1^{c_1}\cdots y_n^{c_n}$ for $(\lambda,\beta)=(\sum_{i=1}^na_i\rho_i,\sum_{i=1}^nc_i\alpha_i)\in P\oplus Q$\\\hline
$(\zeta_i=z^{v_i}:i\in I)$&$(\hy_1,\ldots,\hy_n,x_1,\ldots,x_n)$ for $\hy_i=y_ix_1^{b_{1i}}\cdots x_n^{b_{ni}}$\\\hline
$\zeta^{(\beta,\lambda)}$&
$x^\lambda \hy^\beta=x_1^{a_1}\cdots x_n^{a_n}\hy_1^{c_1}\cdots\hy_n^{c_n}$ for $(\beta,\lambda)=(\sum_{i=1}^nc_i\alpha_i,\sum_{i=1}^na_i\rho_i)\in Q\oplus P$\\\hline
$\k$&a field of characteristic zero\\\hline
$\k\llbracket \zeta\rrbracket $&$\k\llbracket x,\hy\rrbracket =\k\llbracket x_1,\ldots,x_n,\hy_1,\ldots,\hy_n\rrbracket $\\\hline
$\m$&ideal in $\k\llbracket x,\hy\rrbracket $ consisting of series with constant term zero\\\hline
\end{longtable}}% 




We work in the formal power series ring $\k\llbracket x,\hy\rrbracket $ or, sometimes for convenience, in the quotient $\k\llbracket x,\hy\rrbracket /\m^{k+1}$ for $k\ge 0$.
A \newword{wall} $(\d,f_\d)$ consists of a codimension-$1$ cone $\d$ in $V^*$ and a function $f_\d\in\k\llbracket \hy\rrbracket $ (or in $\k\llbracket \hy\rrbracket /\m^{k+1}$) such that:
\begin{enumerate}[label=(\roman*)]%[(i)]
\item $\d$ is contained in $\beta^\perp$ for some primitive $\beta\in Q^+$ and defined by inequalities of the form $\br{p,\phi}\le0$ for $\phi\in Q$.
\item $f_\d=f_\d(\hy^\beta)$ is in the univariate power series ring $\k\llbracket \hy^\beta\rrbracket $ for this primitive $\beta$ (or $f_\d\in\k\llbracket \hy^\beta\rrbracket /\br{\hy^{(k+1)\beta}}$).
\end{enumerate}
A wall $(\d,f_\d)$ is \newword{incoming} if the vector $\omega(\,\cdot\,,\beta)\in V^*$ is in $\d$ and otherwise it is \newword{outgoing}.
Two walls are \newword{parallel} if they are contained in the same hyperplane.

A \newword{scattering diagram} is a collection $\D$ of walls such that the set $\D_k$ of walls $(\d,f_\d)\in\D$ with $f_\d\not\equiv 1$ modulo $\m^{k+1}$ is finite for all $k\ge 1$.
The \newword{support} $\Supp(\D)$ is the union of the walls of $\D$.

Given a scattering diagram $\D$, a \newword{generic} path for $\D$ is a piecewise differentiable path $\gamma:[0,1]\to V^*$ that:
\begin{itemize}
\item does not pass through the intersection of any two non-parallel walls of $\D$;
\item does not pass through the relative boundary of any wall;
\item has endpoints $\gamma(0)$ and $\gamma(1)$ contained in $V^*\setminus\Supp(\D)$; and
\item crosses walls only transversely.
\end{itemize}
%Proposition~\ref{generic exists} says that given $p,q\in V^*\setminus\Supp(\D)$, there exists a generic path $\gamma$ for $\D$ with $\gamma(0)=p$ and $\gamma(1)=q$.

Suppose $\gamma$ is a generic path for $\D$ and $(\d,f_\d)$ is a wall of $\D$ with $\gamma(t)\in\d$ for some $t\in(0,1)$.
The \newword{wall-crossing automorphism} $\p_{\gamma,\d,t}:\k\llbracket x,\hy\rrbracket \to\k\llbracket x,\hat y\rrbracket $ associated to this crossing is given by 
\begin{align}
\label{theta def x}
\p_{\gamma,\d,t}(x^\lambda)&=x^\lambda f_\d^{\br{\lambda,\pm\beta\ck}},\\
\label{theta def hat y}
\p_{\gamma,\d,t}(\hy^\phi)&=\hy^\phi f_\d^{\omega(\pm\beta\ck\!,\,\phi)},
\end{align}
where $\beta\ck$ is the normal vector to $\d$ that is contained in $Q^+$ and is primitive in $Q\ck$, taking $+\beta\ck$ if $\br{\gamma'(t),\beta}<0$ or $-\beta\ck$ if $\br{\gamma'(t),\beta}>0$.
(If $\gamma$ is not differentiable at $t$, the sign of $\br{\gamma'(t),\beta}$ still makes sense, recording the direction in which $\gamma$ crosses the wall.)
The explicit dependence on $t$ in the notation $\p_{\gamma,\d,t}$ is meant to emphasize that $\gamma$ might cross $\d$ multiple times, and in different directions, so that the sign chosen in~\eqref{theta def x} and~\eqref{theta def hat y} depends on more than just $\gamma$ and $\d$.

For each $k\ge 1$, define $\p_{\gamma,\D_k}$ to be $\p_{\gamma,\d_\ell,t_\ell}\circ\cdots\circ\p_{\gamma,\d_1,t_1}:\k\llbracket x,\hy\rrbracket \to\k\llbracket x,\hy\rrbracket $ such that $\d_1,\ldots,\d_\ell$ is the sequence of walls of $\D_k$ crossed by $\gamma$ with $\d_i$ crossed at time $t_i$ and $t_1\le t_2\le\cdots\le t_\ell$.
(There is a finite sequence of crossings because $\D_k$ is finite and because $\gamma$ is generic.)
Define the \newword{path-ordered product} $\p_{\gamma,\D}:\k\llbracket x,\hy\rrbracket \to\k\llbracket x,\hy\rrbracket $ to be $\lim_{k\to\infty}\p_{\gamma,\D_k}$.
We say $\D$ is \newword{consistent} if $\p_{\gamma,\D}$ depends only on $\gamma(0)$ and $\gamma(1)$.
By~\cite[Proposition~2.4]{scatfan}, $\D$ is consistent if and only if each $\D_k$ is consistent modulo $\m^{k+1}$.
When $\D$ is consistent and $p,q\in V^*\setminus\Supp(\D)$, we define $\p_{p,q,\D}=\p_{\gamma,\D}$ for $\gamma$ a generic path from $p$ to $q$, which exists by~\cite[Proposition~2.2]{scatfan}.

It is useful to reinterpret $\p_{\gamma,\D}$ as a map on Laurent monomials sending $x^\lambda y^\phi$ to $x^\lambda y^\phi f_\d^{\br{\lambda,\pm\beta\ck}}$ (with the choice of sign for $\pm\beta\ck$ as in the definition).
With that interpretation, the following is~\cite[Proposition~2.5]{scatfan}, translated into our setup.
\begin{prop}\label{just x's} 
A path-ordered product $\p_{\gamma,\D}$ is determined entirely by its values $\p_{\gamma,\D}(x_1),\ldots,\p_{\gamma,\D}(x_n)$, or equivalently by its values $\p_{\gamma,\D}(x^\lambda)$ for $\lambda\in P$.
\end{prop}

Two scattering diagrams $\D$ and $\D'$ are \newword{equivalent} if and only if $\p_{p,q,\D}=\p_{p,q,\D'}$ for all $p,q\in V^*\setminus(\Supp(\D)\cup\Supp(\D'))$.
A \newword{general} point is a point $p\in V^*$ contained in at most one hyperplane $\beta^\perp$ with $\beta\in Q^+$.
Given a scattering diagram $\D$ and a general point $p$, write $f_p(\D)=\prod_{\d\ni p}f_\d\in\k\llbracket \hy\rrbracket $.
By~\cite[Lemma~1.9]{GHKK}, two scattering diagrams $\D$ and $\D'$ are equivalent if and only $f_p(\D)$ and $f_p(\D')$ agree on all general points $p$.
A scattering diagram has \newword{minimal support} if no equivalent scattering diagram has strictly smaller support.
Every consistent scattering diagram is equivalent to a scattering diagram $\D$ with minimal support and such that each $\D_k$ has minimal support~\cite[Proposition~2.8]{scatfan}.

Given a scattering diagram $\D$ and $\beta\in Q^+$, the \newword{rampart} of $\D$ associated to $\beta$ is the union of all supports of walls of $\D$ contained in $\beta^\perp$. 
%(Typically, the rampart of $\D$ associated to $\beta$ is a smaller set than $\beta^\perp$.)
For $p\in V^*$, write $\Ram_\D(p)$ for the set of ramparts $R$ of $\D$ such that $p\in R$, and write $\D\setminus\Ram_\D(p)$ for the set of walls of $\D$ not contained in any rampart in $\Ram_\D(p)$.

If $\D$ is consistent and has minimal support, we say $p,q\in V^*$ are $\D$-equivalent if and only if there is a path $\gamma$ from $p$ to $q$ on which $\Ram_\D(\,\cdot\,)$ is constant.
This happens if and only if $\Ram_\D(p)=\Ram_\D(q)$ and $p$ and $q$ are in the same path-connected component of $(\cap\Ram_\D(p))\setminus(\Supp(\D\setminus\Ram_\D(p)))$.
The closure of a $\D$-class is a \newword{$\D$-cone}.

A \newword{closed convex cone} in $V^*$ is a subset of $V^*$ that is closed in the usual sense, and also closed under addition and closed under nonnegative scaling.
A subset $F$ of a closed convex cone $C$ is a \newword{face} of $C$ if it is cone and has the property that if $x,y\in C$ and $F$ contains some point in the line segment $\seg{xy}$ besides $x$ and $y$, the whole segment $\seg{xy}$ is in $F$.
A \newword{fan} is a set $\F$ of closed convex cones such that if $C\in\F$ then every face of $C$ is in $\F$, and such that if $C,D\in\F$ then $C\cap D$ is a face of $C$ and a face of $D$.
We write $|\F|$ for the union of the cones in $\F$.
A fan $\F$ is \newword{complete} if $|\F|$ is the entire ambient vector space.

If $\D$ is consistent and has minimal support, then each $\D$-cone is a closed convex cone~\cite[Proposition~3.5]{scatfan} and the collection $\Fan(\D)$ of all $\D$-cones and their faces is a complete fan in $V^*$ by~\cite[Theorem~3.1]{scatfan}.

\subsection{Transposed cluster scattering diagrams}\label{trans clus sec}
An exchange matrix $B$ determines the \newword{transposed cluster scattering diagram with principal coefficients} $\Scat^T (B)$.
This is the unique (up to equivalence) consistent scattering diagram that is obtained by appending \emph{outgoing} walls to an initial scattering diagram $%\set
\{(\alpha_i^\perp,1+ \hy_i): i\break =1,\ldots,n\}$.
In the more general (``non-transposed'') setup of~\cite{scatfan} (as described in Table~\ref{scat root dict}), the scattering diagram $\Scat^T(B)$ is $\Scat(\tB,\s)$ for $\tB=[B^T\,\,I_n]$, where $I_n$ is an identity matrix, and $\s=(\alpha_1,\ldots,\alpha_n,\rho_1,\ldots,\rho_n)$.
We write $\Scat^T(B)$ rather than $\Scat(B^T)$ because we think of $B$ as our primary object, and we want to think of the global transpose not as a modification of $B$, but as a choice of conventions for creating a scattering diagram.
The fan $\Fan(\Scat^T(B))$ is denoted $\ScatFan^T(B)$ and called the \newword{transposed scattering fan}.

%SinceArXiv: Changed the theorem and the short paragraphs before and after to reflect the error that Bridgeland mentioned to me.
We pause here to quote a result that supports the use of root systems in our setup for scattering diagrams.

\begin{theorem}\label{root thm}
If $B$ is skew-symmetric and acyclic, then every wall of $\Scat^T(B)$ is normal to a root.
\end{theorem}

\emph{A priori}, every wall is normal to a positive vector in the root lattice, not necessarily a root.
Theorem~\ref{root thm} follows from~\cite[Proposition~10.1]{GHKK}.
(See also~\cite[Example~10.4]{GHKK}.) 
It also follows from~\cite[Lemma~11.4]{Bridgeland} when $B$ is non-degenerate, relaxing the acyclicity requirement but requiring the existence of a genteel potential. 
We expect that the theorem extends to the case where $B$ is merely skew-symmetrizable.

The following useful fact about transposed cluster scattering diagrams with principal coefficients is proved by applying the antipodal map throughout and observing that the sign changes cancel when we check consistency.
\begin{prop}\label{scat antip}
For any exchange matrix $B$, 
\[\Scat^T(-B)=\set{(-\d,f_\d((\hy')^\beta)):\,(\d,f_\d(\hy^\beta)\in\Scat^T(B)},\]
where the $\hy_i$ are the monomials defined as in Table~\ref{scat root data} using $B$ while the $\hy'_i$ are defined using $-B$, and $f_\d((\hy')^\beta)$ is obtained from $f_\d(\hy^\beta)$ by replacing each $\hy_i$ by $\hy_i'$.
\end{prop}

We write $\A_\bullet(B)$ for the cluster algebra with principal coefficients associated to $B$, in the sense of~\cite[Definition~3.1]{ca4}.
As discussed above, we have passed from \emph{wide} to \emph{tall} extended exchange matrices, and this is the reason for dealing with \emph{transposed} scattering diagrams.
Thus $\A_\bullet(B)$ is the cluster algebra associated (as in~\cite[Definition~2.12]{ca4}) to the extended exchange matrix $\begin{bsmallmatrix}B\\I\end{bsmallmatrix}$, where $I$ is an identity matrix.

A \newword{cluster monomial} is a monomial in the cluster variables in some seed of $\A_\bullet(B)$.
(Conventions on cluster monomials differ, but we take monomials in the ``unfrozen'' cluster variables, not including the ``frozen''/tropical variables.)
We quote two constructions of cluster monomials for $\A_\bullet(B)$.
The first is in terms of broken lines.

Fix a point $p$ in the dominant chamber $D$ such that $p$ is not contained in any hyperplane $\beta^\perp$ for $\beta\in Q$.
We will define a \newword{theta function} $\thet_\lambda$ for every nonzero weight $\lambda\in P\setminus\set{0}$.
Our definition will not depend on the choice of $p$.
(This is not obvious, but rather follows from~\cite[Theorem~3.5]{GHKK}, which is a special case of results of~\cite[Section~4]{CPS}.)

Let $\gamma:(-\infty,0]\to V^*$ be a piecewise linear path with finitely many domains of linearity.
To each domain $L$ of linearity of $\gamma$, assign a monomial $c_Lx^{\lambda_L}y^{\beta_L}$ with $c_L\in\k$ and $(\lambda_L,\beta_L)\in P\oplus Q$.
Then $\gamma$ is a \newword{broken line} for $\lambda$ with endpoint $p$ if the path and the monomials satisfy the following conditions. 
\begin{enumerate}[label=(\roman*)]%[(\roman{enumi})]
\item $\gamma(0)=p$.
\item \label{brok generic root}
$\gamma$ is disjoint from all relative boundaries of walls of $\Scat^T(B)$ and disjoint from all intersections of non-parallel walls of $\Scat^T(B)$. 
\item In each domain $L$ of linearity, $\gamma'$ is constantly equal to $-\lambda_L$.
\item If $L$ is the unbounded domain of linearity of $\gamma$, then $c_Lx^{\lambda_L}y^{\beta_L}=x^\lambda$.
\item \label{change slope}
At each point $t$ of nonlinearity, passing (as the parameter increases) from a domain $L$ of linearity to a domain $L'$ of linearity, by~\ref{brok generic root} there exists $\beta\ck$ primitive in $Q\ck$ such that all walls containing $\gamma(t)$ are in $(\beta\ck)^\perp$ and ${\br{\lambda_L,\beta\ck}>0}$.
If $f$ is the product of the $f_\d$ for all walls $(\d,f_\d)$ with $\gamma(t)\in\d$, then $c_{L'}x^{\lambda_{L'}}y^{\beta_{L'}}$ equals $c_Lx^{\lambda_L}y^{\beta_L}$ times a term in $f^{\br{\lambda_L,\beta\ck}}$.
\end{enumerate}
These conditions in particular allow us to recover the monomials from the path $\gamma$.

Writing $c_\gamma x^{\lambda_\gamma}y^{\beta_\gamma}$ for the monomial on the domain of linearity containing $0$, we define the theta function $\thet_\lambda$ to be the sum, over all broken lines for $\lambda$ with endpoint $p$, of the monomials $c_\gamma x^{\lambda_\gamma}y^{\beta_\gamma}$.
This is an element of $x^\lambda \k\llbracket \hy\rrbracket $.
(We emphasize that the monomial on a domain $L$ of linearity is $c_Lx^{\lambda_L}y^{\beta_L}$, not $c_Lx^{\lambda_L}\hy^{\beta_L}$.
Thus when the monomial changes as described in~\ref{change slope}, both $\lambda_L$ and $\beta_L$ change.)

The subtleties inherent in computing theta functions are compounded by the appearance of both roots and co-roots in the definition.
We give some examples of computing theta functions in rank $2$ in Section~\href{https://alco.centre-mersenne.org/articles/10.5802/alco.107/}{3.4}.

The dominant chamber $D$ is a cone in $\ScatFan^T(B)$. 
Write $\ChamberFan^T(B)$ for the subfan of $\ScatFan^T(B)$ consisting of $D$, all maximal cones $D'$ adjacent to $D$, all maximal cones adjacent to such $D'$, etc., together with all faces of these cones.
The notation $\ChamberFan^T(B)$ will be short-lived in this paper, as almost immediately it will be replaced by a more enlightening notation.
The following is immediate from~\cite[Theorem~5.2]{scatfan} (a version of~\cite[Theorem~4.9]{GHKK}), from~\cite[Corollary~5.9]{GHKK}, and from the definition of $\g$-vectors in~\cite[Section~6]{ca4}.

\begin{theorem}\label{clus mon thm}
The map $\lambda\mapsto\thet_\lambda$ is a bijection from $P\cap|\ChamberFan^T(B)|$ to the set of cluster monomials in $\A_\bullet(B)$.
If $\lambda\in P\cap|\ChamberFan^T(B)|$, then $\lambda$ is the $\g$-vector of the cluster monomial $\thet_\lambda$.
There is a bijection from rays of $\ChamberFan^T(B)$ to cluster variables in $\A_\bullet(B)$ sending each ray to $\thet_\lambda$, where $\lambda$ is the shortest vector in $P$ contained in the ray.
\end{theorem}

The $\g$-vector is defined to be an integer vector, but we interpret elements of~$P$ as $\g$-vectors by taking fundamental-weight coordinates.
Define $\gFan(B)$ to be the set of all cones $C$ such that $C$ is the nonnegative linear span of the $\g$-vectors of a subset of a cluster of $\A_\bullet(B)$.
The following dual version of~\cite[Theorem~0.8]{GHKK} is an immediate corollary of Theorem~\ref{clus mon thm}.

\begin{coro}\label{cham g}
The set $\gFan(B)$ is a fan and coincides with $\ChamberFan^T(B)$.
\end{coro}

Accordingly, we call $\gFan(B)$ the \newword{$\g$-vector fan} of $B$, and we will refer to $\gFan(B)$ rather than $\ChamberFan^T(B)$ through the rest of the paper.

\begin{rema}\label{Clear fr}
In~\cite[Theorem~5.2]{scatfan}, there is an operator $\Clear_{\fr}$ that does not appear in Theorem~\ref{clus mon thm} because the latter concerns principal coefficients.
In the language of~\cite{scatfan}, this is because each term of each $\thet_\lambda$ contains only positive powers of the $\hy_i$, each of which contains only positive powers of the frozen variables~$y_j$.
\end{rema}

The following is~\cite[Theorem~4.6]{scatfan} in our transposed principal-coefficients setting.
\begin{theorem}\label{clus easy root}
If $F$ and $G$ are adjacent maximal cones of $\gFan(B)$, then the function $f_p(\Scat^T(B))$ is $1+\hy^{\beta}$ for every general point $p$ in $F\cap G$, where $\beta$ is the primitive normal to $F\cap G$ in $Q^+$.
\end{theorem}

Our second construction of cluster monomials is in terms of path-ordered products.
The following is a rephrasing of~\cite[Theorem~5.6]{scatfan}.

\begin{theorem}\label{clus mon pop root}
If $\lambda\in P$ is contained in a cone of $\gFan(B)$, then the cluster monomial $\thet_\lambda$ with $\g$-vector $\lambda$ is $\p_{q,p,\Scat^T(B)}(x^\lambda)$ for any point $p$ in the interior of the dominant chamber $D$ and any point $q$ in the interior of a maximal cone $C$ of $\gFan(B)$ such that $\lambda\in C$.
\end{theorem}

The function $\p_{q,p,\Scat^T(B)}(x^\lambda)$ is $F_\lambda\cdot x^\lambda$ for some $F_\lambda\in\k\llbracket \hy\rrbracket $.
By~\cite[Corollary~6.3]{ca4}, the cluster variable with $\g$-vector $\lambda$ is $x^\lambda$ times a polynomial in the $\hy$ called the \newword{$F$-polynomial} of the cluster variable.
(This works because in the principal coefficients case, the denominator in~\cite[(6.5)]{ca4} is $1$.)
Combining Theorem~\ref{clus mon thm} and Theorem~\ref{clus mon pop root}, we have the following immediate corollary.

\begin{coro}\label{clus mon F}
If $\lambda\in P$ is the $\g$-vector of a cluster variable, the corresponding $F$-polynomial is $x^{-\lambda}\cdot\p_{q,p,\Scat^T(B)}(x^\lambda)$ for any point $p$ in the interior of the dominant chamber $D$ and any point $q$ in the interior of a maximal cone $C$ of $\gFan(B)$ such that $\lambda\in C$.
\end{coro}

\section{Scattering diagrams of rank2}\label{rk2 sec}
\subsection{Rank-2 infinite type} \label{rk2 inf sec}
We now consider the transposed cluster scattering diagram with principal coefficients for the exchange matrix $B=\begin{bsmallmatrix}0&b\\a&0\end{bsmallmatrix}$ with $ab\le-4$.
Up to the symmetry of swapping the indices $1$ and $2$, we may as well assume that $a<0$, so that $b>0$.
Continuing the notation above, we have $n=2$ and we have initial cluster variables $x_1$ and $x_2$.
For each $i\in\integers$, we define $x_{i+1}$ to be the cluster variable in $\A_\bullet(B)$ obtained by exchanging $x_{i-1}$ out of the cluster $\set{x_{i-1},x_i}$.
Thus, equivalently, $x_{i-1}$ is obtained by exchanging $x_{i+1}$ out of the cluster $\set{x_i,x_{i+1}}$.
Write $\g_i$ for the $\g$-vector of $x_i$.
In our pictures and verbal descriptions, we will let $\rho_1$ point to the right of the plane and $\rho_2$ point up, give them the same length \emph{in the page}, and talk about slopes of lines in the usual sense for this placement of $\rho_1$ and $\rho_2$.

We follow~\cite[Section~9]{universal} in characterizing the $\g_i$ directly. (Compare~\cite[Example~1.15]{GHKK} and~\cite[Section~3]{CGMMRSW}.)
We begin by defining a polynomial $P_m$ in $ab$ for each $m\ge-2$.
Set $P_{-2}=-1$ and $P_{-1}=0$, and for $m\ge0$, define
\begin{equation}\label{Pm recur}
P_m=\begin{cases}
-abP_{m-1}-P_{m-2}&\text{if }m\text{ is even, or}\\
P_{m-1}-P_{m-2}&\text{if }m\text{ is odd.}\\
\end{cases}
\end{equation}
Several of the $P_m$ are shown in Table~\ref{pm table}.

\begin{table}[ht]
%\begin{tabular}{|r||l|}
%\hline
%$m$&$P_m$\\\hline\hline
%$-2$&	$-1$	\\\hline
%$-1$&	$0$	\\\hline
%$0$&	$1$	\\\hline
%$1$&	$1$	\\\hline
%$2$&	$-ab-1$\\\hline
%$3$&	$-ab-2$\\\hline
%$4$&	$a^2b^2+3ab+1$	\\\hline
%%$5$&	$a^2b^2+4ab+3$	\\\hline
%%$6$&	$-a^3b^3-5a^2b^2-6ab-1$	\\\hline
%%$7$&$-a^3b^3-6a^2b^2-10ab-4$		\\\hline
%%$8$&$a^4b^4+7a^3b^3+15a^2b^2+10ab+1$		\\\hline
%%$9$&$a^4b^4+8a^3b^3+21a^2b^2+20ab+5$		\\\hline
%%$10$&$-a^5b^5-9a^4b^4-28a^3b^3-35a^2b^2-15ab-1$		\\\hline
%\end{tabular}\\[3pt]
\renewcommand{\arraystretch}{1.2}
\begin{tabular}{|r|||c|c|c|c|c|c|c|c|}
\hline
$m$&$-2$&$-1$&$0$&$1$&$2$&$3$&$4$&$5$\\\hline\hline
$P_m$&	$-1$	&$0$	&$1$	&$1$	&$-ab-1$&$-ab-2$&$a^2b^2+3ab+1$	&$a^2b^2+4ab+3$	\\\hline
%$6$&	$-a^3b^3-5a^2b^2-6ab-1$	\\\hline
%$7$&$-a^3b^3-6a^2b^2-10ab-4$		\\\hline
%$8$&$a^4b^4+7a^3b^3+15a^2b^2+10ab+1$		\\\hline
%$9$&$a^4b^4+8a^3b^3+21a^2b^2+20ab+5$		\\\hline
%$10$&$-a^5b^5-9a^4b^4-28a^3b^3-35a^2b^2-15ab-1$		\\\hline
\end{tabular}\\[3pt]
\caption{The polynomials $P_m$}
\label{pm table}
\end{table}

The polynomials $P_m$ are defined in~\cite[Section~9]{universal} using a summation formula, and a specific relationship~\cite[(9.10)]{universal} is described between the $P_m$ and the Chebyshev polynomials of the second kind.
The recursive definition given here for the $P_m$ follows by the defining recursion for the Chebyshev polynomials.
As observed in~\cite[Section~9]{universal}, each $P_m$ is positive for $m\ge0$.
By~\cite[Proposition~9.6]{universal}, the $\g$-vectors are given by
\begin{equation}\label{ba gs}
\g_i=\begin{cases}
-P_{-i-1}\rho_1-aP_{-i-2}\rho_2&\text{if }i\text{ is odd and }i\le-1,\\
-bP_{-i-1}\rho_1+P_{-i-2}\rho_2&\text{if }i\text{ is even and }i\le0,\\
-P_{i-3}\rho_1-aP_{i-2}\rho_2&\text{if }i\text{ is odd and }i\ge1,\\
-bP_{i-3}\rho_1+P_{i-2}\rho_2&\text{if }i\text{ is even and }i\ge2,\\
\end{cases}
\end{equation}
Write $C_i$ for the cone spanned by the $\g$-vectors of $x_i$ and $x_{i+1}$.
In particular, $C_1$ is the dominant chamber $D$.
We write $\cc_i$ for the positive root orthogonal to $\g_i$. 
Recalling that the fundamental roots $\rho_i$ are dual to the simple \emph{co-roots} $\alpha_i\ck=\delta_i^{-1}\alpha_i$ and observing that the diagonalizing factors in this case must be 
$\delta_1=\frac{\gcd(-a,b)}{b}$ and $\delta_2=\frac{\gcd(-a,b)}{-a}$, we compute
%\begin{equation}\label{ba ccks}
%\c\ck_i=\begin{cases}
%P_{-i-2}\alpha\ck_1+bP_{-i-1}\alpha\ck_2&\text{if }i\text{ is even and }i\le-2,\\
%-aP_{-i-2}\alpha\ck_1+P_{-i-1}\alpha\ck_2&\text{if }i\text{ is odd and }i\le-1,\\
%\alpha\ck_1&\text{if }i=0,\\
%\alpha\ck_2&\text{if }i=1,\text{ or},\\
%P_{i-2}\alpha\ck_1+bP_{i-3}\alpha\ck_2&\text{if }i\text{ is even and }i\ge2,\\
%-aP_{i-2}\alpha\ck_1+P_{i-3}\alpha\ck_2&\text{if }i\text{ is odd and }i\ge3,\\
%\end{cases}
%\end{equation}
%\begin{equation}\label{ba ccks in root basis} 
%\c\ck_i=\begin{cases}
%P_{-i-2}\frac{b}{\gcd(-a,b)}\alpha_1+bP_{-i-1}\frac{-a}{\gcd(-a,b)}\alpha_2&\text{if }i\text{ is even and }i\le-2,\\
%-aP_{-i-2}\frac{b}{\gcd(-a,b)}\alpha_1+P_{-i-1}\frac{-a}{\gcd(-a,b)}\alpha_2&\text{if }i\text{ is odd and }i\le-1,\\
%\frac{b}{\gcd(-a,b)}\alpha_1&\text{if }i=0,\\
%\frac{-a}{\gcd(-a,b)}\alpha_2&\text{if }i=1,\text{ or},\\
%P_{i-2}\frac{b}{\gcd(-a,b)}\alpha_1+bP_{i-3}\frac{-a}{\gcd(-a,b)}\alpha_2&\text{if }i\text{ is even and }i\ge2,\\
%-aP_{i-2}\frac{b}{\gcd(-a,b)}\alpha_1+P_{i-3}\frac{-a}{\gcd(-a,b)}\alpha_2&\text{if }i\text{ is odd and }i\ge3,\\
%\end{cases}
%\end{equation}
\begin{equation}\label{ba cs} 
\cc_i=\begin{cases}
P_{-i-2}\alpha_1-aP_{-i-1}\alpha_2&\text{if }i\text{ is even and }i\le-2,\\
bP_{-i-2}\alpha_1+P_{-i-1}\alpha_2&\text{if }i\text{ is odd and }i\le-1,\\
\alpha_1&\text{if }i=0,\\
\alpha_2&\text{if }i=1,\text{ or},\\
P_{i-2}\alpha_1-aP_{i-3}\alpha_2&\text{if }i\text{ is even and }i\ge2,\\
bP_{i-2}\alpha_1+P_{i-3}\alpha_2&\text{if }i\text{ is odd and }i\ge3,\\
\end{cases}
\end{equation}

The fan $\gFan(B)$ covers all of $V^*$ except the positive linear span $C_\infty$ of the vectors 
\begin{align}
\begin{split}
\g_\infty&=-2\sqrt{-ab}\,\rho_1-a(\sqrt{-ab}+\sqrt{-ab-4})\,\rho_2\quad\text{and}\\
\g_{-\infty}&=-b(\sqrt{-ab}+\sqrt{-ab-4})\,\rho_1+2\sqrt{-ab}\,\rho_2.
\end{split}
\end{align}
These vectors are in the limiting directions of $\g_i$ as $i\to\infty$ and $i\to-\infty$ respectively.
As we discuss in Section~\ref{affrk2 sec}, in the affine case (when $ab=-4$), the vectors $\g_\infty$ and $\g_{-\infty}$ are parallel, so that $C_\infty$ is a single limiting ray.


By Theorem~\ref{clus mon thm} and Theorem~\ref{clus easy root}, the walls of $\Scat^T(B)$ not contained in $C_\infty$ are the rays spanned by the $\g_i$, each marked with the function $(1+\hy^{\cc_i})$.
The remaining possibilities for walls are the rational rays contained in $C_\infty$.
We consider a scattering diagram in the equivalence class of $\Scat^T(B)$ with one wall in each such rational ray, without assuming that the attached functions are nontrivial.
Thus, for each rational number in reduced form $p/q$ with $p>0$ and $q<0$ and with 
\begin{equation}\label{pq range}
\frac{a(\sqrt{-ab}+\sqrt{-ab-4})}{2\sqrt{-ab}}\le\frac{\,p\,}q\le\frac{2\sqrt{-ab}}{-b(\sqrt{-ab}+\sqrt{-ab-4})}\,,
\end{equation}
let $(\d_{p/q},f_{p/q})$ be the wall such that $\d_{p/q}$ is the ray spanned by $q\rho_1+p\rho_2$.
Then $f_{p/q}$ is a formal power series in $\hy^{\cc_{p/q}}$, where $\cc_{p/q}=\frac{bp}{\gcd(bp,aq)}\alpha_1+\frac{aq}{\gcd(bp,aq)}\alpha_2$ is the primitive element of $Q^+$ that is orthogonal to $q\rho_1+p\rho_2$.
The primitive element of $Q\ck$ parallel to $\cc_{p/q}$ is $\cc\ck_{p/q}=p\alpha\ck_1-q\alpha\ck_2$.

Fix $p/q$ satisfying~\eqref{pq range} and choose a path $\gamma_\infty:(0,1]\to V^*\setminus\set{0}$ and a path ${\gamma_{-\infty}:(0,1]\to V^*\setminus\set{0}}$ such that
\begin{itemize} 
\item $\lim_{t\to0^+}\gamma_{\pm\infty}(t)=\g_{\pm\infty}$
\item $\gamma_{\pm\infty}(1)=\rho_1+\rho_2$
\item $\gamma_\infty$ moves in a strictly monotone-clockwise manner about $0$,
%the measure of the angle $\angle(\gamma_\infty(t),0,\rho_1+\rho_2)$ is strictly decreasing.
\item $\gamma_{-\infty}$ moves in a strictly monotone-counterclockwise manner about $0$.
%the measure of the angle $\angle(\gamma_{-\infty}(t),0,\rho_1+\rho_2)$ is strictly increasing.
\end{itemize}
%Thus $\gamma_\infty$ moves clockwise about $0$, while $\gamma_{-\infty}$ moves counterclockwise.
For each $i\ge1$, let $\gamma_i$ be a subpath of $\gamma_{\infty}$ starting in the interior of $C_i$ and ending at $\rho_1+\rho_2$.
For $i\le1$, let $\gamma_i$ be a subpath of $\gamma_{-\infty}$ starting in the interior of $C_i$ and ending at $\rho_1+\rho_2$.
For consistency, take $\gamma_1$ to be the constant path at $\rho_1+\rho_2$.

We abbreviate $\p_{\gamma,\Scat^T(B)}$ as $\p_\gamma$ for any path $\gamma$.
We define $\p_\infty=\lim_{i\to\infty}\p_{\gamma_i}$.
This limit exists because for all $k$, there exists an $i_k$ such that, for $i\ge i_k$, the path obtained by deleting $\gamma_i$ from $\gamma_\infty$ crosses no wall of $\Scat^T(B)_k$.
We also define $\p_{-\infty}=\lim_{i\to-\infty}\p_{\gamma_i}$, and this limit exists for the analogous reason.
Indeed, the limit definitions are only necessary when $ab=-4$. 
When $ab<-4$, we may as well complete $\gamma_\infty$ and $\gamma_{-\infty}$ to paths from the full interval $[0,1]$ to $V^*\setminus\set{0}$. 
These paths are generic because $\g_\infty$ and $\g_{-\infty}$ are not contained in a rational ray, and the limits $\p_\infty$ and $\p_{-\infty}$ coincide with $\p_{\gamma_\infty}$ and $\p_{\gamma_{-\infty}}$.

We also define paths $\gamma_\pm:(0,1]\to V^*\setminus\set{0}$ such that
\begin{itemize} 
\item $\lim_{t\to0^+}\gamma_{\pm\infty}(t)=q\rho_1+p\rho_2$
\item $\gamma_{\pm\infty}(1)=\g_{\pm\infty}$
%\item the measure of the angle $\angle(\gamma_+(t),0,\rho_1+\rho_2)$ is strictly decreasing.
%\item the measure of the angle $\angle(\gamma_-(t),0,\rho_1+\rho_2)$ is strictly increasing.
\item $\gamma_+$ moves in a strictly monotone-clockwise manner about $0$,
\item $\gamma_-$ moves in a strictly monotone-counterclockwise manner about $0$.
\end{itemize}
Write $\p_\pm$ for the path-ordered products obtained from the paths $\gamma_\pm$, taking appropriate limits as above.
When $ab=-4$, we think of $\gamma_\pm$ as empty paths and treat an automorphism $\p_\pm$ as the identity when it appears in formulas.

Write $\p_{p/q}$ for the wall-crossing automorphism $\p_{\gamma,\d_{p/q}}$ for a path $\gamma$ that crosses $\d_{p/q}$ with derivative $\rho_1+\rho_2$.
%Consistency of the scattering diagram says that ${\p_{-\infty}\circ\p_-\circ\p_{p/q}^{-1}\circ\p_+^{-1}\circ\p_\infty^{-1}}$ is the identity.
%This consistency allows us to make the following general statement about $f_{p/q}$.

Since $p>0$ and $q<0$, in particular $p-q>0$.
Since $f_{p/q}$ is a (univariate) formal power series in $\hy^{\cc_{p/q}}$, the following proposition determines $f_{p/q}$ completely.

\begin{prop}\label{diagonal terms ba}
For all $k\ge0$, the coefficient of $\hy^{k\cc_{p/q}}$ in $f_{p/q}$ equals the coefficient of $\hy^{k\cc_{p/q}}$ in $\bigl(\frac{x_1x_2}{\p_{-\infty}(\p_-(x_1x_2))}\bigr)^{\frac1{p-q}}$.
Each $\hy^{k\cc_{p/q}}$ is $\hy_1^i\hy_2^j$ with $\frac ji=\frac{aq}{bp}$, and every other term in $\bigl(\frac{x_1x_2}{\p_{-\infty}(\p_-(x_1x_2))}\bigr)^{\frac1{p-q}}$ involves $\hy_1^i\hy_2^j$ with $\frac ji<\frac{aq}{bp}$.
\end{prop}


\begin{proof}
Consistency of the scattering diagram says that ${\p_{-\infty}\circ\p_-\circ\p_{p/q}^{-1}\circ\p_+^{-1}\circ\p_\infty^{-1}}$ is the identity map, so we apply it to $\p_\infty(\p_+(\p_{p/q}(x_1x_2))$ to obtain
\begin{equation}\label{diag terms ba starting point}
\p_{-\infty}(\p_-(x_1x_2))=\p_\infty(\p_+(\p_{p/q}(x_1x_2)).
\end{equation}
We calculate $\p_{p/q}(x_1x_2)=x_1x_2 f_\d^{\br{\rho_1+\rho_2,-p\alpha_1\ck+q\alpha_2\ck}}=x_1x_2f_{p/q}^{q-p}$.
Thus the quantity $\p_\infty(\p_+(\p_{p/q}(x_1x_2)))$ is computed by starting with $x_1x_2f_{p/q}^{q-p}$ and repeatedly replacing a monomial by the same monomial times an integer power of a power series in $\hy^\beta$ for various $\beta\in Q^+$.
Specifically, each $\beta$ is of the form $c_1\alpha_1+c_2\alpha_2$ such that the slope $\frac{c_2}{c_1}$ of the ray spanned by $\beta$ is positive but strictly less than $\frac{aq}{pb}$, which is the slope of the ray spanned by $\cc_{p/q}$.

Therefore, by~\eqref{diag terms ba starting point}, every term in $\frac{\p_{-\infty}(\p_-(x_1x_2))}{x_1x_2}$ involves $\hy_1^i\hy_2^j$ with $\frac ji\le\frac{aq}{bp}$, and the terms where $\frac ji=\frac{aq}{bp}$ are exactly $f_{p/q}^{q-p}$.
Thus, to find $f_{p/q}$, we raise $\frac{\p_{-\infty}(\p_-(x_1x_2))}{x_1x_2}$ to the power $\frac1{q-p}$ and restrict to terms involving $\hy_1^i\hy_2^j$ with $\frac ji=\frac{aq}{bp}$.
\end{proof}

By definition, $\p_{-\infty}(\p_-(x_1x_2))$ is $\p_{-\infty}(x_1x_2\cdot\M_{p/q}(\hy_1,\hy_2))$ for some $\M_{p/q}(\hy_1,\hy_2)$ in $\k\llbracket \hy\rrbracket $.
Thus 
\begin{equation}\label{p M}
\p_{-\infty}(\p_-(x_1x_2))=\p_{-\infty}(x_1x_2)\cdot\p_{-\infty}(\M_{p/q}(\hy_1,\hy_2)).
\end{equation}
The factor $\p_{-\infty}(\M_{p/q}(\hy_1,\hy_2))$ is difficult to deal with in general.
This is the reason we eventually restrict to the affine case, where $\p_{-\infty}(\M_{p/q}(\hy_1,\hy_2))=1$.
However, before restricting to the affine case, we make a general observation about the factor $\p_{-\infty}(x_1x_2)$.
The observation uses the following lemma, which is easily verified using~\eqref{Pm recur} and~\eqref{ba gs}.
\begin{lemma}\label{deal with gs}
For $i\le-2$, \,
$\g_i=\begin{cases}
b\g_{i+1}-\g_{i+2}&\text{if }i\text{ is even, or}\\
-a\g_{i+1}-\g_{i+2}&\text{if }i\text{ is odd.}\\
\end{cases}$
\end{lemma}
%\begin{proof}
%If $i$ is even, then
%\begin{align*}
%b\g_{i+1}-\g_{i+2}&=b(-P_{-i-2}\rho_1-aP_{-i-3}\rho_2)-(-bP_{-i-3}\rho_1+P_{-i-4}\rho_2)\\
%&=-b(P_{-i-2}-P_{-i-3})\rho_1+(-abP_{-i-3}-P_{-i-4})\rho_2\\
%&=-bP_{i-1}\rho_1+P_{-i-2}\rho_2\\
%&=\g_i
%\end{align*}
%If $i$ is odd, then
%\begin{align*}
%-a\g_{i+1}-\g_{i+2}&=-a(-bP_{-i-2}\rho_1+P_{-i-3}\rho_2)-(-P_{-i-3}\rho_1-aP_{-i-4}\rho_2)\\
%&=-(-abP_{-i-2}-P_{-i-3})\rho_1-a(P_{-i-3}-P_{-i-4})\rho_2\\
%&=-P_{i-1}\rho_1-aP_{-i-2}\rho_2\\
%&=\g_i
%\end{align*}
%\end{proof}


We will also need a polynomial $Q_m$ in $a$ and $b$ for $m\ge0$, given by $Q_0=1$ and $Q_1=-1$ and by the recursion 
\begin{equation}\label{Qm recur}
Q_m=\begin{cases}
bQ_{m-1}-Q_{m-2}&\text{if }m\text{ is even, or}\\
-aQ_{m-1}-Q_{m-2}&\text{if }m\text{ is odd.}
\end{cases}
\end{equation}
Several values of $Q_m$ are shown in Table~\ref{qm table}.

\begin{table}[ht]
\renewcommand{\arraystretch}{1.2}
%\begin{tabular}{|r||l|}
%\hline
%$m$&$Q_m$\\\hline\hline
%$0$&	$1$	\\\hline
%$1$&	$-1$	\\\hline
%$2$&	$-b-1$\\\hline
%$3$&	$ab+a+1$\\\hline
%$4$&	$ab^2+ab+2b+1$	\\\hline
%%$5$&	$-a^2b^2-a^2b-3ab-2a-1$	\\\hline
%%$6$&	$-a^2b^3-a^2b^2-4ab^2-3ab-3b-1$	\\\hline
%%$7$&	$a^3b^3+a^3b^2+5a^2b^2+4a^2b+6ab+3a+1$		\\\hline
%\end{tabular}\\[3pt]
\begin{tabular}{|r||c|c|c|c|c|c|c|c|}
\hline
$m$&$0$&$1$&$2$&$3$&$4$&$5$\\\hline\hline
$Q_m$&	{\small$1$}	&{\small$-1$}	&{\small$-b-1$}&{\small$ab+a+1$}&{\small$ab^2+ab+2b+1$}&{\small$-a^2b^2-a^2b-3ab-2a-1$}	\\\hline
%$6$&	$-a^2b^3-a^2b^2-4ab^2-3ab-3b-1$	\\\hline
%$7$&	$a^3b^3+a^3b^2+5a^2b^2+4a^2b+6ab+3a+1$		\\\hline
\end{tabular} \\[3pt]
\caption{The polynomials $Q_m$}
\label{qm table}
\end{table}

\begin{prop}\label{p inf lim}
If $F_i$ is the $F$-polynomial of the cluster variable $x_i$, then
\begin{equation}\label{p inf lim eq}
\frac{x_1x_2}{\p_{-\infty}(x_1x_2)}=\lim_{i\to-\infty}F_i^{Q_{-i-1}}\cdot F_{i+1}^{-Q_{-i}}.
\end{equation}
\end{prop}
\begin{proof}
We use Theorem~\ref{clus mon pop root} and Theorem~\ref{clus mon thm} to write $x_i=\p_{\gamma_i}(x^{\g_i})$ and $x_{i+1}=\p_{\gamma_i}(x^{\g_{i+1}})$ for all $i\in\integers$.
We are justified in using the same path $\gamma_i$ in both of these equations because $C_i$ is the cone spanned by $\g_i$ and $\g_{i+1}$.
By an easy induction using Lemma~\ref{deal with gs} and~\eqref{Qm recur}, we show that $-\g_iQ_{-i-1}+\g_{i+1}Q_{-i}=\rho_1+\rho_2$ for all $i\le-1$.
%For $i=-1$, this says $-\rho_2\cdot(-1)-(-\rho_1)\cdot1=\rho_1+\rho_2$.
%For $i=-2$, it says $(-\rho_1)\cdot(-b-1)-(-b\rho_1+\rho_2)\cdot(-1)=\rho_1+\rho_2$.
%For $i<-2$ even, 
%\begin{align*}
%\g_{i+1}Q_{-i}-\g_iQ_{-i-1}&=\g_{i+1}(bQ_{-i-1}-Q_{-i-2})-(b\g_{i+1}-\g_{i+2})Q_{-i-1}\\
%&=-\g_{i+1}Q_{-i-2}+\g_{i+2}Q_{-i-1}\\
%&=\rho_1+\rho_2
%\end{align*}
%For $i<-2$ odd, 
%\begin{align*}
%\g_{i+1}Q_{-i}-\g_iQ_{-i-1}&=\g_{i+1}(-aQ_{-i-1}-Q_{-i-2})-(-a\g_{i+1}-\g_{i+2})Q_{-i-1}\\
%&=-\g_{i+1}Q_{-i-2}+\g_{i+2}Q_{-i-1}\\
%&=\rho_1+\rho_2
%\end{align*}
Thus $x_i^{-Q_{-i-1}}x_{i+1}^{Q_{-i}}=\p_{\gamma_i}(x_1x_2)$.
Thus $\p_{-\infty}(x_1x_2)$ is $\lim_{i\to-\infty}x_i^{-Q_{-i-1}}x_{i+1}^{Q_{-i}}$, which equals $x_1x_2\lim_{i\to-\infty}F_i^{-Q_{-i-1}}\cdot F_{i+1}^{Q_{-i}}$.
\end{proof}

\subsection{Rank-2 affine type}\label{affrk2 sec}
In the affine cases (when $ab=-4$), the considerations of Section~\ref{rk2 inf sec} are sufficient to determine the function attached to the limiting wall.
We will prove the following theorem.
The remaining affine cases can be obtained from the theorem by Proposition~\ref{scat antip} and/or by swapping the indices $1$ and $2$.


\begin{thm}\label{rk2 aff formula}
The function on the limiting wall of $\Scat^T\bigl(\begin{bsmallmatrix}\,\,\,\,\,0&\,2\\-2&\,0\end{bsmallmatrix}\bigr)$ is $\displaystyle\frac1{(1-\hy_1\hy_2)^2}$.
The function on the limiting wall of $\Scat^T\bigl(\begin{bsmallmatrix}\,\,\,\,\,0&\,1\\-4&\,0\end{bsmallmatrix}\bigr)$ is $\displaystyle\frac{1+\hy_1\hy_2^2}{(1-\hy_1\hy_2^2)^2}$.
\end{thm}

We start with the first assertion, which is the case $a=-2$ and $b=2$ in the notation of Section~\ref{rk2 inf sec}.
For these values of $a$ and $b$, the slope of the limiting ray is $p/q=-1$, the positive root orthogonal to that ray is $\cc_{-1}=\alpha_1+\alpha_2$, and $Q_m=-2m+1$.
Combining Propositions~\ref{diagonal terms ba} and~\ref{p inf lim} in this case, we obtain the following proposition.
\begin{prop}\label{diagonal terms}
If $a=-2$ and $b=2$, then for all $k\ge0$, the coefficient of $\hy_1^k\hy_2^k$ in $f_{-1}$ equals the coefficient of $\hy_1^k\hy_2^k$ in $\sqrt{\lim_{i\to-\infty}F_i^{2i+3}\cdot F_{i+1}^{-2i-1}}$.
\end{prop}

We will call a term \newword{diagonal} if it is a constant times $\hy_1^k\hy_2^k$ and \newword{super-diagonal} if it is a constant times $\hy_1^j\hy_2^k$ for $k>j$.
The following lemma will let us evaluate the limit in Proposition~\ref{diagonal terms}.

\begin{lemma}\label{F coeffs}
When $a=-2$ and $b=2$, the $F$-polynomial $F_i$ has the following properties for $i\le 0$.
\begin{enumerate}[label=(\roman*)]%[\rm(i)]
\item \label{leading}
As a polynomial in $\hy_1$ (with coefficients polynomials in $\hy_2$), the leading term is $\hy_1^{-i}(1+\hy_2)^{-i+1}$.
\item \label{others}
The only super-diagonal term is $\hy_1^{-i}\hy_2^{-i+1}$.
\item \label{key terms}
The diagonal terms are $(j+1)\hy_1^j\hy_2^j$ for $0\le j\le-i$.
\end{enumerate}
\end{lemma}
\begin{proof}
We argue by induction on $-i$.
Unwrapping~\cite[Proposition~5.1]{ca4} in this case, we obtain $F_0=1+\hy_2$ and $F_{-1}=1+\hy_1(1+\hy_2)^2$, and for $i\le-2$ we obtain
\begin{equation}\label{F recur}
F_i=\frac{F_{i+1}^2+\hy_1^{-i-2}\hy_2^{-i-1}}{F_{i+2}}.
\end{equation}
(The exponents on $\hy_1$ and $\hy_2$ come from~\eqref{ba cs}.)
We have established the base cases $i=0$ and $i=-1$, and for the rest of the proof, we take $i<-1$.
By induction and by dividing leading terms, assertion~\ref{leading} follows easily from~\eqref{F recur}.

To prove the remaining assertions, we track the diagonal and super-diagonal terms through~\eqref{F recur}.
We perform long division and, since we know that the result is a polynomial, we can ignore terms that are not diagonal or super-diagonal.
Indeed, because by induction we know that the highest order term in the denominator of~\eqref{F recur} is super-diagonal, we need only track super-diagonal terms in the numerator.
By induction, the relevant terms in the numerator are 
\[\hy_1^{-2i-2}\hy_2^{-2i}+\hy_1^{-i-2}\hy_2^{-i-1}+\hy_1^{-i-1}\hy_2^{-i}\sum_{j=0}^{-i-1}2(j+1)\hy_1^j\hy_2^j\]
and the relevant terms in the denominator are $\hy_1^{-i-2}\hy_2^{-i-1}+\sum_{j=0}^{-i-2}(j+1)\hy_1^j\hy_2^j$.
Division yields super-diagonal and diagonal terms $\hy_1^{-i}\hy_2^{-i+1}+\sum_{j=0}^{-i}(j+1)\hy_1^j\hy_2^j$ as desired.
\end{proof}

Lemma~\ref{F coeffs}\ref{key terms} implies that, for $i<0$, the restriction of $F_i$ to diagonal terms agrees with the formal power series $\sum_{k\ge0}(k+1)\hy_1^k\hy_2^k=(1-\hy_1\hy_2)^{-2}$ up to the term $\hy_1^{-i}\hy_2^{-i}$.
By Lemma~\ref{F coeffs}\ref{others}, the only super-diagonal terms in $F_i$ have degree at least $-i$ in $\hy_1$, so the diagonal terms of $F_{i}^{2i+3}\cdot F_{i+1}^{-2i-1}$ agree, up to $\hy_1^{-i-1}\hy_2^{-i-1}$, with $(1-\hy_1\hy_2)^{-4}$.
Thus Proposition~\ref{diagonal terms} establishes the first assertion of Theorem~\ref{rk2 aff formula}.

\begin{rema}\label{UMN REU}
In~\cite{Zhang}, there is a formal power series closely related to the function $f_{-1}$ on the limiting wall.
It appears as a limit of ``stable cluster variables'' (certain transformed $F$-polynomials). 
\end{rema}

We next prove the second assertion of Theorem~\ref{rk2 aff formula}.
In this case, $a=-4$ and $b=1$, the slope of the limiting ray is $p/q=-2$, the positive root orthogonal to that ray is $\cc_{-2}=\alpha_1+2\alpha_2$, and $Q_m$ is $-\frac32m+1$ for $m$ even or $-3m+2$ for $m$ odd.
We know that the limit in Proposition~\ref{p inf lim} exists, so we are free to approach it on the even values of $i$.
We will do this by replacing $i$ by $2i$ in the limit formula.
Thus Propositions~\ref{diagonal terms ba} and~\ref{p inf lim} give the following proposition.
\begin{prop}\label{diagonal terms asym}
If $a=-4$ and $b=1$, then for all $k\ge0$, the coefficient of $\hy_1^k\hy_2^{2k}$ in $f_{-2}$ equals the coefficient of $\hy_1^k\hy_2^{2k}$ in $\sqrt[3]{\lim_{i\to-\infty}F_{2i}^{6i+5}\cdot F_{2i+1}^{-3i-1}}$.
%\[\frac{x_1x_2}{\p_{-\infty}(x_1x_2)}=\lim_{i\to-\infty}F_{2i}^{Q_{-2i-1}}\cdot F_{2i+1}^{-Q_{-2i}}\]
%\[\frac{x_1x_2}{\p_{-\infty}(x_1x_2)}=\lim_{i\to-\infty}F_{2i}^{-3(-2i-1)+2}\cdot F_{2i+1}^{-(-\frac32(-2i)+1)}\]
%\[\frac{x_1x_2}{\p_{-\infty}(x_1x_2)}=\lim_{i\to-\infty}F_{2i}^{(6i+3)+2}\cdot F_{2i+1}^{-(3i+1)}\]
\end{prop}

Again, evaluating the limit will require determining some relevant coefficients of the $F$-polynomials, but this time we must separate the even and odd indices.
We will also \emph{re-use} the words diagonal and super-diagonal to deal with this case.
We call a term \newword{diagonal} if it is a constant times $\hy_1^k\hy_2^{2k}$ and \newword{super-diagonal} if it is a constant times $\hy_1^k\hy_2^\ell$ for $\ell>2k$.
Furthermore, a term is \newword{super-super-diagonal} if it is a constant times $\hy_1^k\hy_2^\ell$ for $\ell>2k+1$.

\begin{lemma}\label{F coeffs asym}
Suppose $a=-4$ and $b=1$.
The even-indexed $F$-polynomials $F_{2i}$ for $i\le0$ have the following properties:
\begin{enumerate}[label=(\roman*)]%[\rm(i)]
\item \label{leading asym}
As a polynomial in $\hy_1$ (with coefficients polynomials in $\hy_2$), the leading term is $\hy_1^{-i}(1+\hy_2)^{-2i+1}$.
\item \label{super asym}
The only super-diagonal term is $\hy_1^{-i}\hy_2^{-2i+1}$.
\item \label{key terms asym}
The diagonal terms are $\sum_{j=0}^{-i}(2j+1)\hy_1^j\hy_2^{2j}$.
%For $0\le j\le-i$, the coefficient of $\hy_1^j\hy_2^{2j}$ is $2j+1$.
\end{enumerate}
The odd-indexed $F$-polynomials $F_{2i+1}$ for $i\le-1$ have the following properties:
\begin{enumerate}[label=(\roman*)] %%[\rm(i)] 
\setcounter{enumi}{3}
\item \label{leading asym odd}
As a polynomial in $\hy_1$ (with coefficients polynomials in $\hy_2$), the leading term is $\hy_1^{-2i-1}(1+\hy_2)^{-4i}$.
\item \label{super asym odd}
The only super-diagonal terms are $\hy_1^{-2i-1}\hy_2^{-4i}+\sum_{j=1}^{-i}4j\hy_1^{-i+j-1}\hy_2^{-2i+2j-1}$.
\item \label{key terms asym odd}
%The coefficients of the form $\hy_1^k\hy_2^{2k}$ are $\bigl(\sum_{j=0}^{-i}(2j+1)\hy_1^j\hy_2^{2j}\bigr)^2+\sum_{j=0}^{-i}b_{ij}\hy_1^{-i+j}\hy_2^{-2i+2j}$.
The diagonal terms agree with $\bigl(\sum_{j\ge0}(2j+1)\hy_1^j\hy_2^{2j}\bigr)^2$ up to $(-2i+1)\hy_1^{-i}\hy_2^{-2i}$.
%For $0\le j\le-i$, the coefficient of $\hy_1^j\hy_2^{2j}$ is $\displaystyle\frac{(j+1)(2j^2+4j+3)}3$ and for $-i+1\le j\le-2i-1$, the coefficient of $\hy_1^j\hy_2^{2j}$ is $\displaystyle\frac{WHAT?}{}$
\end{enumerate}
\end{lemma}

\begin{proof}
We argue by induction on $m\le0$ that $F_m$ has the properties described in the proposition.
In this case~\cite[Proposition~5.1]{ca4} says that $F_0=1+\hy_2$ and ${F_{-1}=1+\hy_1(1+\hy_2)^4}$.
Also, for $i\le -1$
\begin{equation}\label{F recur asym even}
F_{2i}=\frac{F_{2i+1}+\hy_1^{-i-1}\hy_2^{-2i-1}}{F_{2i+2}}.
\end{equation}
and for $i\le-2$,
\begin{equation}\label{F recur asym odd}
F_{2i+1}=\frac{(F_{2i+2})^4+\hy_1^{-2i-3}\hy_2^{-4i-4}}{F_{2i+3}}.
\end{equation}
(The monomials $\hy_1^{-i-1}\hy_2^{-2i-1}$ and $\hy_1^{-2i-3}\hy_2^{-4i-4}$ are $\hy^{\cc_{m-1}}$ for $m=2i$ or $2i+1$ and $\cc_{m-1}$ as in~\eqref{ba cs}.)
We have established the base cases $m=0$ and $m=-1$, and for the rest of the proof, we take $m<-1$.
Assertions~\ref{leading asym} and~\ref{leading asym odd} follow easily by induction.

To prove the remaining assertions, we track diagonal and super-diagonal terms through the right sides of~\eqref{F recur asym even} and~\eqref{F recur asym odd}.
Since we know that the right side is a polynomial, we can simply perform polynomial long division and ignore all terms that are not diagonal or super-diagonal.
Because we know by induction the highest-order terms of the denominator of~\eqref{F recur asym even}, we need only concern ourselves with super-diagonal terms in the numerator.
Similarly, in the numerator of~\eqref{F recur asym odd}, we need only concern ourselves with super-super-diagonal terms.

For $m=2i$, by induction the relevant terms in the numerator in~\eqref{F recur asym even} are 
\begin{equation}\label{even numer}
\hy_1^{-2i-1}\hy_2^{-4i}+\sum_{j=1}^{-i}4j\hy_1^{-i+j-1}\hy_2^{-2i+2j-1}+\hy_1^{-i-1}\hy_2^{-2i-1}
\end{equation}
and the relevant terms in the denominator are
\begin{equation}\label{even denom}
\hy_1^{-i-1}\hy_2^{-2i-1}+\sum_{j=0}^{-i-1}(2j+1)\hy_1^j\hy_2^{2j}.
\end{equation}
The relevant terms of the quotient are $\hy_1^{-i}\hy_2^{-2i+1}+\sum_{j=0}^{-i}(2j+1)\hy_1^j\hy_2^{2j}$, as desired.

For $m=2i+1$, by induction, the relevant terms in the numerator are the super-super-diagonal terms in $\bigl(\hy_1^{-i-1}\hy_2^{-2i-1}+\sum_{j=0}^{-i-1}(2j+1)\hy_1^j\hy_2^{2j}\bigr)^4+\hy_1^{-2i-3}\hy_2^{-4i-4}$, namely
\begin{multline}\label{odd numer}
\hy_1^{-4i-4}\hy_2^{-8i-4}+4\hy_1^{-3i-3}\hy_2^{-6i-3}\sum_{j=0}^{-i-1}(2j+1)\hy_1^j\hy_2^{2j}\\
+6\hy_1^{-2i-2}\hy_2^{-4i-2}\biggl(\sum_{j=0}^{-i-1}(2j+1)\hy_1^j\hy_2^{2j}\biggr)^2
+\hy_1^{-2i-3}\hy_2^{-4i-4}.
\end{multline}
The relevant terms in the denominator are
\begin{multline}\label{odd denom}
\hy_1^{-2i-3}\hy_2^{-4i-4}+\hy_1^{-i-2}\hy_2^{-2i-3}\sum_{j=0}^{-i-1}4j\hy_1^j\hy_2^{2j}\\
+\biggl(\sum_{j=0}^{-i-1}(2j+1)\hy_1^j\hy_2^{2j}\biggr)^2+\hy_1^{-i}\hy_2^{-2i}\sum_{j=0}^{-i-3}b_j\hy_1^j\hy_2^{2j},
\end{multline}
for some coefficients $b_j$ that (it will turn out) are not important.
The super-diagonal terms in the quotient are $\hy_1^{-2i-1}\hy_2^{-4i}+4\hy_1^{-i-1}\hy_2^{-2i-1}\sum_{j=0}^{-i}j\hy_1^j\hy_2^{2j}$ and the diagonal terms are
\begin{multline}\label{X}
1+(6\hy_1\hy_2^2-\hy_1^2\hy_2^4)\biggl(\sum_{j=0}^{-i-1}(2j+1)\hy_1^j\hy_2^{2j}\biggr)^2\\-16\biggl(\sum_{j=0}^{-i-1}j\hy_1^j\hy_2^{2j}\biggr)\biggl(\sum_{j=0}^{-i}j\hy_1^j\hy_2^{2j}\biggr)+\hy_1^{-i+2}\hy_2^{-i+4}\biggl(\sum_{j=0}^{-i-1}b_j\hy_1^j\hy_2^{2j}\biggr).
\end{multline}
But~\eqref{X} agrees, up to terms involving $\hy_1^{-i}\hy_2^{-2i}$ with 
\begin{equation}\label{X FPS}
1+(6\hy_1\hy_2^2-\hy_1^2\hy_2^4)\biggl(\frac{1+\hy_1\hy_2^2}{(1-\hy_1\hy_2^2)^2}\biggr)^2\\-16\biggl(\frac{\hy_1\hy_2^2}{(1-\hy_1\hy_2^2)^2}\biggr)=\biggl(\frac{1+\hy_1\hy_2^2}{(1-\hy_1\hy_2^2)^2}\biggr)^2,
\end{equation}
which equals $=\bigl(\sum_{j\ge0}(2j+1)\hy_1^j\hy_2^{2j}\bigr)^2$ as desired.
\end{proof}

Lemma~\ref{F coeffs asym} says in particular that for $i\le-1$ the diagonal terms of $F_{2i}$ agree, up to $\hy_1^{-i}\hy_2^{-2i}$, with $\frac{1+\hy_1\hy_2^2}{(1-\hy_1\hy_2^2)^2}$ and the diagonal terms of $F_{2i+1}$ agree, up to $\hy_1^{-i}\hy_2^{-2i}$, with $\Bigl(\frac{1+\hy_1\hy_2^2}{(1-\hy_1\hy_2^2)^2}\Bigr)^2$.
Furthermore, the only super-diagonal terms in either polynomial have degree at least $-i$ in $\hy_1$, so the diagonal terms of $F_{2i}^{6i+5}\cdot F_{2i+1}^{-3i-1}$ agree, up to $\hy_1^{-i-1}\hy_2^{-2i-2}$, with $\Bigl(\frac{1+\hy_1\hy_2^2}{(1-\hy_1\hy_2^2)^2}\Bigr)^{6i+5}\cdot\Bigl(\frac{1+\hy_1\hy_2^2}{(1-\hy_1\hy_2^2)^2}\Bigr)^{-6i-2}$, which equals $\Bigl(\frac{1+\hy_1\hy_2^2}{(1-\hy_1\hy_2^2)^2}\Bigr)^3$.
Now Proposition~\ref{diagonal terms asym} completes the proof of Theorem~\ref{rk2 aff formula}.

\begin{thebibliography}{99}
\bibitem[3]{Bridgeland} Tom Bridgeland, ``Scattering diagrams, Hall algebras and stability conditions'', \emph{Algebr. Geom.}\textbf{4}(2017),no.5,523--561.
\bibitem[5]{CPS} Michael Carl, Max Pumperla and Bernd Siebert, \emph{A tropical view of Landau-Ginzburg models}, \url{http://www.math.uni-hamburg.de/home/siebert/preprints/LGtrop.pdf},2010.
\bibitem[6]{CGMMRSW} Man Wai Cheung, Mark Gross, Greg Muller, Gregg Musiker, Dylan Rupel, Salvatore Stella and Harold Williams, ``The greedy basis equals the theta basis: a rank two haiku'', \emph{J. Comb. Theory,Ser.A}\textbf{145}(2017),150--171.
\bibitem[8]{ca4} Sergey Fomin and Andrei Zelevinsky, ``Cluster algebras.IV.Coefficients'', \emph{Compos. Math.}\textbf{143}(2007),no.1,112--164.
\bibitem[9]{GHKK} Mark Gross, Paul Hacking, Sean Keel and Maxim Kontsevich, ``Canonical bases for cluster algebras'', \emph{J. Am. Math. Soc.}\textbf{31}(2018),no.2,497--608.
\bibitem[24]{universal} Nathan Reading, ``Universal geometric cluster algebras'', \emph{Math. Z.}\textbf{277}(2014),no.1-2,499--547.
\bibitem[25]{scatfan} Nathan Reading, ``Scattering Fans'', \emph{Int. Math. Res. Not.}(2018),rny260.
\bibitem[34]{Zhang} Grace Zhang, \emph{Stable cluster variables}, REU Report, \url{http://www-users.math.umn.edu/~reiner/REU/Zhang2016.pdf},2016.
\end{thebibliography}
\end{document}
